% !TeX root = ../DoABook.tex

\chapter{نمای کلی الگوریتم‌های پایه تخمین \lr{DOA}}
\label{chap:3}

\section{مقدمه}
\label{sec:3-1}

این فصل نمایی کلی از برخی الگوریتم‌های محبوب برای تخمین \lr{DOA} را ارائه می‌دهد. به‌طور کلی، تکنیک‌های تخمین جهت ورود (\lr{DOA}) را می‌توان به‌طور گسترده به سه دسته تکنیک‌های شکل‌دهی پرتو متعارف\LTRfootnote{Conventional beamforming techniques}، تکنیک‌های مبتنی بر زیرفضا\LTRfootnote{Subspace-based techniques} و تکنیک‌های بیشینه درست‌نمایی\LTRfootnote{Maximum likelihood techniques} تقسیم‌بندی کرد.

ساختار این فصل به شرح زیر است: در بخش \ref{sec:3-1}، مدل سیگنال و داده مورد استفاده در این کتاب توصیف می‌شود. برای تشریح این مدل، از یک آرایه خطی یکنواخت استفاده شده است. در بخش \ref{sec:3-2}، مفهوم آرایه‌های حسگری مرکز-متقارن\LTRfootnote{Centro-symmetric sensor arrays} که مورد نیاز بسیاری از الگوریتم‌های \lr{DOA} هستند، مورد بحث قرار می‌گیرد. دو نوع از این آرایه‌ها، یعنی آرایه خطی یکنواخت (\lr{ULA})\LTRfootnote{Uniform Linear Array (ULA)} و آرایه مستطیلی یکنواخت (\lr{URA})\LTRfootnote{Uniform Rectangular Array (URA)}، توصیف می‌شوند. در ادامه، اصول اولیه چند الگوریتم پایه متعلق به هر یک از دسته‌های \lr{DOA} به‌طور خلاصه مورد بحث قرار می‌گیرند.

یک رساله دکتری عالی در زمینه پردازش سیگنال آرایه‌ای برای تخمین‌های \lr{DOA} توسط \lr{M. Haardt} در انتشارات \textit{شاکر ورلاگ} منتشر شده است \cite{1}. این اثر یک بررسی و مطالعه جامع روی یکی از تکنیک‌های \lr{DOA}، یعنی تکنیک \lr{ESPRIT}، ارائه می‌دهد. کتاب حاضر از آن به‌عنوان یکی از مراجع اصلی استفاده می‌کند. به‌منظور تسهیل در بررسی این مرجع اصلی توسط خواننده و جلوگیری از سردرگمی، این کتاب از همان نمادهای ریاضی و اصطلاحات فنی به کار رفته در \cite{1} استفاده می‌کند؛ این نمادها و اصطلاحات در سایر متون علمی مرتبط نیز به کار گرفته شده‌اند.


\section{مدل داده}
\label{sec:3-2}

اکثر رویکردهای مدرن در پردازش سیگنال، مدل‌مبنا\LTRfootnote{Model-based} هستند؛ به این معنی که بر فرضیات خاصی در مورد داده‌های مشاهده‌شده در دنیای واقعی تکیه دارند. مدل داده‌ی غالب که در ادامه این کتاب مورد استفاده قرار می‌گیرد، در این بخش بحث می‌شود.

در سرتاسر توصیف ما از الگوریتم‌های تخمین \lr{DOA}، سناریوی زیر فرض و حفظ شده است \cite{2}:

\begin{itemize}
	\item \textbf{محیط انتشار همسانگرد و خطی\LTRfootnote{Isotropic and linear transmission medium}:} تعداد $d$ منبع وجود دارند که $d$ سیگنال را گسیل می‌کنند؛ این $d$ سیگنال از طریق یک محیط عبور کرده و سپس بر یک آرایه آنتنی یا حسگری $M$-عنصری می‌تابند. فرض می‌شود محیط انتشار بین منابع و آرایه، همسانگرد و خطی است؛ به عبارت دیگر، محیط دارای ویژگی‌های فیزیکی یکسان در تمامی جهت‌های مختلف است و سیگنال‌ها یا امواج در هر نقطه خاص را می‌توان به‌صورت خطی برهم‌نهی کرد. ویژگی همسانگرد و خطی بودن محیط تضمین می‌کند که: (۱) ویژگی انتشار امواج با تغییر \lr{DOA} سیگنال‌ها تغییر نمی‌کند، و (۲) سیگنال‌هایی که در محیط حرکت کرده و سپس توسط هر عنصر از آرایه $M$-عنصری دریافت می‌شوند، می‌توانند به‌عنوان برهم‌نهی خطی $d$ جبهه موج سیگنالِ تولیدشده توسط $d$ منبع محاسبه شوند. علاوه بر این، بهره هر آنتن یا عنصر حسگر برابر با یک فرض می‌شود.
	
	\item \textbf{فرض میدان‌دور\LTRfootnote{Far-field assumption}:} منابع سیگنال $d$ در فاصله‌ای دور از آرایه قرار دارند، به‌طوری که جبهه موج تولیدشده توسط هر منبع با جهت انتشار یکسان به تمامی عناصر می‌رسد و جبهه موج تخت است (تقریب میدان‌دور). بنابراین، میدان‌های در حال انتشار $d$ سیگنال که به آرایه رسیده‌اند، نسبت به یکدیگر موازی در نظر گرفته می‌شوند. این فرض به‌طور کلی با بزرگتر در نظر گرفتن فاصله بین منابع سیگنال و آرایه نسبت به ابعاد آرایه آنتن محقق می‌شود. یک قاعده سرانگشتی این است که فاصله بزرگتر از $2D^2/\lambda$ باشد که در آن $D$ ابعاد آرایه و $\lambda$ طول موج سیگنال‌ها است \cite{3}.
	
	\item \textbf{فرض باندباریک\LTRfootnote{Narrowband assumption}:} سیگنال‌های $d$ از منابع $d$ دارای فرکانس حامل یکسانی هستند، به‌طوری که محتوای فرکانسی آن‌ها در مجاورت فرکانس حامل $f_c$ متمرکز شده است. در این صورت، از نظر ریاضی، هر یک از سیگنال‌های لحظه‌ای که از منابع خارج می‌شوند را می‌توان به‌صورت زیر بیان کرد:
	\begin{equation}
		s_i^r(t) = \alpha_i(t) \cos[2\pi f_c t + \beta_i(t)], \quad 1 \leq i \leq d
		\label{eq:3-1}
	\end{equation}
	سیگنال‌ها تا زمانی باندباریک محسوب می‌شوند که دامنه $\alpha_i(t)$ و فاز حامل اطلاعات $\beta_i(t)$ آن‌ها نسبت به $\tau$ (زمان انتشار سیگنال‌های موج از یک عنصر به عنصر دیگر) به‌کندی تغییر کنند. به عبارت دیگر:
	\begin{equation}
		\alpha_i(t - \tau) \approx \alpha_i(t) \quad \text{و} \quad \beta_i(t - \tau) \approx \beta_i(t)
		\label{eq:3-2}
	\end{equation}
	تغییرات کند $\alpha_i(t)$ و فازهای $\beta_i(t)$ تضمین می‌کنند که تبدیل فوریه رابطه \eqref{eq:3-1} بیشترین محتوای فرکانسی را در همسایگی فرکانس حامل $f_c$ داشته باشد. با تعریف پوش مختلط\LTRfootnote{Complex envelope} سیگنال یا فازور\LTRfootnote{Phasor} به‌صورت زیر، می‌توان به عبارتی مناسب برای تحلیل‌های ریاضی دست یافت:
	\begin{equation}
		s_i(t) = s_{env, i}(t) = \alpha_i(t) e^{j\beta_i(t)} \quad \text{به‌طوری که} \quad s_i^r(t) = \text{Re}\{s_i(t) e^{j 2\pi f_c t}\}
		\label{eq:3-3}
	\end{equation}
	این شکل از سیگنال‌های مختلط (یا تحلیلی) توسط اکثر گیرنده‌ها پشتیبانی می‌شود که سیگنال‌های دریافتی را به هر دو مؤلفه هم‌فاز (بخش حقیقی) و مؤلفه متعامد (بخش موهومی) تجزیه می‌کنند.
	
	\item \textbf{کانال \lr{AWGN}\LTRfootnote{Additive White Gaussian Noise (AWGN)}:} فرض می‌شود نویز یک فرآیند گاوسی سفید مختلط است. نویز جمع‌شونده از یک فرآیند تصادفی با میانگین صفر و از نظر مکانی ناهمبسته\LTRfootnote{Spatially uncorrelated} در نظر گرفته می‌شود که با سیگنال‌ها نیز ناهمبسته است. نویزها در تمامی عناصر آرایه دارای واریانس مشترک $\sigma_N^2$ بوده و بین تمامی عناصر یا حسگرها ناهمبسته هستند.
\end{itemize}

شکل \ref{fig:3-1} سناریوی مربوط به فرضیات فوق را ترسیم می‌کند.


\begin{figure}[ht]
	\centering
	\caption{سناریوی مورد بررسی در این کتاب.}
	\label{fig:3-1}
\end{figure}

\subsection{آرایه خطی یکنواخت (\lr{ULA})}
\label{subsec:3-2-1}

اکنون یک آرایه خطی یکنواخت\LTRfootnote{Uniform Linear Array (ULA)} شامل $M$ عنصر یکسان و همه‌جهته را در نظر بگیرید که روی یک خط تراز شده و دارای فواصل یکسان هستند. فرض کنید فاصله بین دو عنصر مجاور $\Delta$ و فاصله بین منبع و اولین عنصر (راست‌ترین عنصر) $d_d$ باشد. این پیکربندی در شکل \ref{fig:3-2} نشان داده شده است.

فرض کنید یک سیگنال موج تخت تولیدشده توسط منبع $i$ با زاویه $\theta_i$ بر آرایه می‌تابد. سیگنال تولیدشده توسط منبع $i$ یک سیگنال باندباریک $s_i^r(t)$ است که مسافت $d_d$ را با سرعت $c$ طی کرده \cite{3} و به اولین عنصر (راست‌ترین عنصر) می‌رسد. به عبارت دیگر، سیگنال دریافت شده توسط راست‌ترین عنصر، نسخه تأخیریافته‌ی $s_i^r(t)$ با تأخیر $\tau_d = d_d/c$ است. یعنی:
\begin{equation}
	s_{i1}^r(t) = s_i^r(t-\tau_d) = \alpha_i(t-\tau_d) \cos[2\pi f_c(t-\tau_d) + \beta_i(t-\tau_d)] = \text{Re}\{s_{i1}(t) e^{j2\pi f_c t}\}
	\label{eq:3-4}
\end{equation}
که در آن سیگنال فازور متناظر به‌صورت $s_{i1}(t) = \alpha_i(t-\tau_d) e^{j[-2\pi f_c \tau_d + \beta_i(t-\tau_d)]}$ است.

از آنجا که در \lr{ULA} عناصر آرایه روی یک خط قرار دارند، سیگنالی که به عنصر $m$-ام می‌رسد، مسافت بیشتری را نسبت به سیگنال رسیده به راست‌ترین عنصر طی خواهد کرد (شکل \ref{fig:3-2} را ببینید). این مسافت اضافی به‌صورت زیر محاسبه می‌شود:
\begin{equation}
	\Delta_{mi} = (m-1) \Delta \sin \theta_i, \quad m=1, 2, \dots, M
	\label{eq:3-5}
\end{equation}
بنابراین، سیگنالی که به عنصر $m$-ام می‌رسد، دچار تأخیر اضافی می‌شود که برابر است با:
\begin{equation}
	\tau_{mi} = \frac{\Delta_{mi}}{c} = \frac{(m-1) \Delta \sin \theta_i}{c}
	\label{eq:3-6}
\end{equation}

سیگنال دریافت شده توسط عنصر $m$-ام، نسخه تأخیریافته‌ی سیگنال $s_{i1}^r(t)$ (که توسط عنصر اول دریافت شده) با تأخیر اضافی $\tau_{mi}$ است:
\begin{align}
	s_{im}^r(t) &= s_{i1}^r(t-\tau_{mi}) = s_i^r(t-\tau_d-\tau_{mi}) \nonumber \\
	&= \alpha_i(t-\tau_d-\tau_{mi}) \cos[2\pi f_c(t-\tau_d-\tau_{mi}) + \beta_i(t-\tau_d-\tau_{mi})] \nonumber \\
	&\approx \alpha_i(t-\tau_d) \cos[2\pi f_c(t-\tau_d) + \beta_i(t-\tau_d) - (m-1)\mu_i] \nonumber \\
	&= \text{Re}\{s_{i1}(t) e^{j2\pi f_c t} e^{-j(m-1)\mu_i}\}
	\label{eq:3-7}
\end{align}
که در آن $\mu_i = \frac{2\pi f_c \Delta}{c} \sin \theta_i = \frac{2\pi \Delta}{\lambda} \sin \theta_i$ است؛ این مقدار فرکانس مکانی\LTRfootnote{Spatial frequency} نامیده می‌شود که به منبع $i$-ام با زاویه تابش $\theta_i$ اختصاص دارد. $\lambda = c/f_c$ نشان‌دهنده طول موج متناظر با فرکانس حامل $f_c$ است. در استخراج رابطه \eqref{eq:3-7}، تقریب \eqref{eq:3-2} اعمال شده است.

در قالب فازور مختلط، سیگنال دریافتی فوق متناظر است با:
\begin{equation}
	s_{im}(t) \approx \alpha_i(t-\tau_d) e^{j[-2\pi f_c \tau_d + \beta_i(t-\tau_d)]} e^{-j(m-1)\mu_i} = s_{i1}(t) e^{-j(m-1)\mu_i}
	\label{eq:3-8}
\end{equation}

رابطه \eqref{eq:3-8} نشان می‌دهد سیگنالی که توسط عنصر $m$-ام از منبع $i$-ام دریافت می‌شود، مشابه سیگنال دریافت شده توسط عنصر اول (راست‌ترین عنصر) است، اما یک فاکتور جابجایی فاز اضافی به‌صورت $e^{-j(m-1)\mu_i}$ دارد؛ این فاکتور تنها به فرکانس مکانی $\mu_i$ و موقعیت عنصر نسبت به عنصر اول بستگی دارد. برای هر زاویه تابش $\theta_i$ که یک منبع را مشخص می‌کند، یک فرکانس مکانی متناظر $\mu_i$ وجود دارد. بنابراین، هدف کلی از تخمین \lr{DOA} (یعنی $\theta_i$) استخراج این فرکانس مکانی $\mu_i$ از سیگنال‌های دریافت شده توسط آرایه است.

به‌منظور تعیین یکتای $\theta_i$ از روی $\mu_i$، وجود یک تناظر یک‌به‌یک بین آن‌ها الزامی است. در نتیجه، فرکانس‌های مکانی $\mu_i$ به بازه $[-\pi, \pi]$ محدود شده و محدوده \lr{DOA}های ممکن به بازه $[-90^\circ, 90^\circ]$ محدود می‌گردد. این امر به نوبه خود مستلزم آن است که فاصله بین عناصر در شرط $\Delta \leq \lambda/2$ صدق کند. اگر فاصله عناصر در این رابطه صدق نکند، در تعیین \lr{DOA} ابهام به وجود می‌آید، یعنی برای یک مقدار مشخص $\mu_i$، دو راه حل برای زوایا وجود خواهد داشت؛ این وضعیت منجر به ایجاد لوب‌های گریتینگ\LTRfootnote{Grating lobes} در آرایه می‌شود: لوب‌هایی غیر از لوب اصلی که اجازه ورود سیگنال‌ها را از جهت‌های نامطلوب می‌دهند. از نظر فنی، این پدیده علی‌اسینگ مکانی\LTRfootnote{Spatial aliasing} نامیده می‌شود. این مفهوم مشابه نرخ نمونه‌برداری نایکوئیست در تحلیل حوزه فرکانس یک سیگنال است.



\begin{figure}[ht]
	\centering
	\caption{مدل داده برای تخمین \lr{DOA} مربوط به $d$ منبع با یک آرایه خطی $M$-عنصری.}
	\label{fig:3-2}
\end{figure}

اکنون با در نظر گرفتن تمامی سیگنال‌های تولیدشده توسط هر $d$ منبع، $s_i(t)$ برای $1 \le i \le d$، سیگنال کل و نویزهای دریافت شده توسط عنصر $m$-ام در زمان $t$ را می‌توان به‌صورت زیر بیان کرد:
\begin{equation}
	x_m(t) = \sum_{i=1}^{d} s_{im}(t) + n_m(t) = \sum_{i=1}^{d} s_i(t) e^{-j(m-1)\mu_i} + n_m(t), \quad m = 1, 2, \dots, M
	\label{eq:3-9}
\end{equation}


برای تفکیک سیگنال‌های خالص تولیدشده توسط منابع و سیگنال‌های دریافتی یا آشکارسازی‌شده‌ی آلوده به نویز، از این پس به مورد دوم «داده»\LTRfootnote{Data} اطلاق شده و با نماد $\bm{x}$ نشان داده می‌شود.

رابطه \eqref{eq:3-9} را می‌توان به فرم ماتریسی زیر نوشت:
\begin{equation}
	\bm{x}(t) = [\bm{a}(\mu_1), \bm{a}(\mu_2), \dots, \bm{a}(\mu_d)] \bm{s}(t) + \bm{n}(t) = \bm{A}\bm{s}(t) + \bm{n}(t)
	\label{eq:3-10}
\end{equation}
که در آن $\bm{x}(t) = [x_1(t), x_2(t), \dots, x_M(t)]^T$ بردار ستونی داده‌های دریافتی توسط آرایه، $\bm{s}(t) = [s_1(t), s_2(t), \dots, s_d(t)]^T$ بردار ستونی سیگنال‌های تولیدشده توسط منابع و $\bm{n}(t) = [n_1(t), n_2(t), \dots, n_M(t)]^T$ نویزهای جمع‌شونده با میانگین صفر و از نظر مکانی ناهمبسته با ماتریس کوواریانس مکانی برابر با $\sigma_N^2 \bm{I}_M$ است. بردار ستونی هدایت آرایه $\bm{a}(\mu_i)$ به‌صورت زیر تعریف می‌شود:
\begin{equation}
	\bm{a}(\mu_i) = [1, e^{-j\mu_i}, e^{-j2\mu_i}, \dots, e^{-j(M-1)\mu_i}]^T
	\label{eq:3-11}
\end{equation}
این بردار تابعی از فرکانس‌های مکانی ناشناخته $\mu_i$ است و ستون‌های ماتریس هدایت $\bm{A}$ با ابعاد $M \times d$ را تشکیل می‌دهد:
\begin{equation}
	\bm{A} = [\bm{a}(\mu_1), \bm{a}(\mu_2), \dots, \bm{a}(\mu_d)] = 
	\begin{bmatrix} 
		1 & 1 & \dots & 1 \\
		e^{-j\mu_1} & e^{-j\mu_2} & \dots & e^{-j\mu_d} \\
		\vdots & \vdots & \ddots & \vdots \\
		e^{-j(M-1)\mu_1} & e^{-j(M-1)\mu_2} & \dots & e^{-j(M-1)\mu_d}
	\end{bmatrix}
	\label{eq:3-12}
\end{equation}

بسته به پیکربندی‌های مختلف یک آرایه، ماتریس‌های هدایت آرایه‌ای متفاوتی می‌تواند تشکیل شود. بسیاری از الگوریتم‌ها به‌طور خاص مستلزم آن هستند که آرایه‌ها دارای پیکربندی‌های مرکز-متقارن باشند [4]. خوشبختانه، بسیاری از پیکربندی‌های مورد استفاده، مانند آرایه خطی یکنواخت و آرایه مستطیلی، این شرط را برآورده می‌کنند. در بخش بعدی، دو آرایه مرکز-متقارن بسیار رایج را مرور می‌کنیم: \lr{ULA} و \lr{URA}.





\section{آرایه‌های حسگری مرکز-متقارن}
\label{sec:3-3}

در این بخش، تعاریف اولیه آرایه‌های مرکز-هرمیتی\LTRfootnote{Centro-Hermitian} و مرکز-متقارن\LTRfootnote{Centro-symmetric} ارائه می‌شود. این آرایه‌ها به‌دلیل دارا بودن ویژگی‌های خاص، اغلب در کاربردهای عملی مورد استفاده قرار می‌گیرند و معمولاً توسط ماتریس‌های مرکز-متقارن و مرکز-هرمیتی توصیف می‌شوند. در این بخش، بر اساس توصیفات ارائه شده در \cite{1, 5}، به تشریح آن‌ها خواهیم پرداخت.

ابتدا، اجازه دهید $\mathbb{C}^{p \times q}$ را به‌عنوان مجموعه یا فضای ماتریس‌های $p \times q$ در نظر بگیریم. تعاریف زیر بیانگر ماهیت یک ماتریس مرکز-متقارن یا مرکز-هرمیتی است.

\textbf{تعریف ۳-۱} \\
یک ماتریس مختلط $\bm{M} \in \mathbb{C}^{p \times q}$ مرکز-متقارن نامیده می‌شود \cite{1, 5} اگر:
\begin{equation}
	\bm{\Pi}_p \bm{M} \bm{\Pi}_q = \bm{M}
	\label{eq:3-13}
\end{equation}
در اینجا $\bm{\Pi}_p$ ماتریس تعویض\LTRfootnote{Exchange matrix} $p \times p$ است که در پیوست A.2 تعریف شده است؛ این ماتریس روی قطر فرعی (\textit{antidiagonal}) خود دارای مقدار یک و در سایر جاها دارای مقدار صفر است:
\begin{equation}
	\bm{\Pi}_p = 
	\begin{bmatrix}
		0 & \dots & 0 & 1 \\
		0 & \dots & 1 & 0 \\
		\vdots & \dots & \vdots & \vdots \\
		1 & \dots & 0 & 0
	\end{bmatrix}
	\label{eq:3-14}
\end{equation}
اثبات این مطلب دشوار نیست که پیش‌ضرب یک ماتریس در $\bm{\Pi}_p$، ترتیب سطرهای آن ماتریس را معکوس می‌کند (یعنی سطر اول به سطر آخر، سطر دوم به سطر یکی مانده به آخر و سطر آخر به سطر اول تبدیل می‌شود و به همین ترتیب). پس‌ضرب یک ماتریس در $\bm{\Pi}_p$ نیز ترتیب ستون‌های آن را معکوس می‌نماید (یعنی ستون اول به ستون آخر، ستون دوم به ستون یکی مانده به آخر و ستون آخر به ستون اول تبدیل می‌شود). به همین ترتیب، $\bm{\Pi}_q$ ماتریس تعویض $q \times q$ است.

شرط \eqref{eq:3-13} اساساً بیان می‌کند که $\bm{M} \in \mathbb{C}^{p \times q}$ دارای یک مرکز وارونگی\LTRfootnote{Inversion center} است: زمانی که $\bm{M}$ به اندازه ۱۸۰ درجه چرخانده شود، با ماتریس اصلی یکسان خواهد بود. از نظر مکانی، این بدان معناست که به ازای هر نقطه $(x, y, z)$، یک نقطه متمایزناپذیر\LTRfootnote{Indistinguishable} $(-x, -y, -z)$ وجود دارد.

یک ماتریس مختلط $\bm{M} \in \mathbb{C}^{p \times q}$ را مرکز-هرمیتی-متقارن\LTRfootnote{Centro-Hermitian-symmetric} می‌نامند اگر \cite{1, 5}:
\begin{equation}
	\bm{\Pi}_p \bar{\bm{M}} \bm{\Pi}_q = \bm{M}
	\label{eq:3-15}
\end{equation}
که در آن $\bar{\bm{M}}$ مزدوج مختلط $\bm{M}$ است.

شرط \eqref{eq:3-15} در اصل بیان می‌کند که $\bm{M} \in \mathbb{C}^{p \times q}$ دارای یک مرکز وارونگی مزدوج است؛ بدین معنا که وقتی $\bm{M}$ مزدوج شده و سپس ۱۸۰ درجه چرخانده شود، با ماتریس اصلی یکسان می‌گردد. از نظر مکانی، می‌توان دید که به ازای هر نقطه $(x, y, z)$، یک نقطه مزدوج متمایزناپذیر $(-x, -y, -z)$ وجود دارد.

\textbf{تعریف ۳-۲} \\
ماتریس $\bm{Q} \in \mathbb{C}^{p \times q}$ که در رابطه زیر صدق کند:
\begin{equation}
	\bm{\Pi}_p \bar{\bm{Q}} = \bm{Q} \iff \bar{\bm{Q}} = \bm{\Pi}_p \bm{Q}
	\label{eq:3-16}
\end{equation}
ماتریس متقارن حقیقی-چپ-$\bm{\Pi}$\LTRfootnote{Left $\Pi$ real symmetric matrix} نامیده می‌شود.

برای مثال، ماتریس‌های یکانی\LTRfootnote{Unitary matrices}:
\begin{equation}
	\bm{Q}_{2n} = \frac{1}{\sqrt{2}} 
	\begin{bmatrix}
		\bm{I}_n & j\bm{I}_n \\
		\bm{\Pi}_n & -j\bm{\Pi}_n
	\end{bmatrix}
	, \quad 
	\bm{Q}_{2n+1} = \frac{1}{\sqrt{2}} 
	\begin{bmatrix}
		\bm{I}_n & \bm{0} & j\bm{I}_n \\
		\bm{0}^T & \sqrt{2} & \bm{0}^T \\
		\bm{\Pi}_n & \bm{0} & -j\bm{\Pi}_n
	\end{bmatrix}
	\label{eq:3-17}
\end{equation}
به‌ترتیب ماتریس‌های حقیقی-چپ-$\bm{\Pi}$ از مرتبه‌های زوج و فرد هستند. در اینجا $\bm{I}_n$ ماتریس واحد با ابعاد $n \times n$ است. ماتریس‌های حقیقی-چپ-$\bm{\Pi}$ را می‌توان با پس‌ضرب یک ماتریس حقیقی-چپ-$\bm{\Pi}$ مانند $\bm{Q} \in \mathbb{C}^{p \times q}$ در یک ماتریس حقیقی دلخواه $\bm{R}$ نیز به‌دست آورد. یعنی برای هر ماتریس حقیقی-چپ-$\bm{\Pi}$ مانند $\bm{Q}$:
\begin{equation}
	\bm{Q}_r = \bm{QR}, \quad \bm{R} \in \mathbb{R}^{q \times r}
	\label{eq:3-18}
\end{equation}
نیز یک ماتریس حقیقی-چپ-$\bm{\Pi}$ خواهد بود.




\begin{definition}
	یک آرایه حسگری شامل $M$ عنصر، آرایه‌ای مرکز-متقارن\LTRfootnote{Centro-symmetric array} محسوب می‌شود، مشروط بر آنکه مکان عناصر آن نسبت به مرکز هندسی\LTRfootnote{Centroid} متقارن بوده و ویژگی‌های تابشی\LTRfootnote{Radiation characteristics} عناصر جفت‌شده یکسان باشد. بنابراین، می‌توان اثبات کرد که ماتریس هدایت آرایه آن‌ها، $\bm{A}$، در رابطه زیر صدق می‌کند:
	\begin{equation}
		\bm{\Pi}_M \bar{\bm{A}} = \bm{A} \bm{\Lambda}
		\label{eq:3-19}
	\end{equation}
	که در آن $\bm{\Lambda}$ یک ماتریس قطری یکانی\LTRfootnote{Unitary diagonal matrix} است. توجه داشته باشید که در اینجا فرض شده است عناصر حسگر دارای ویژگی‌های تابشی یکسانی هستند.
\end{definition}

در بخش بعدی، دو آرایه مرکز-متقارن پرکاربرد، یعنی آرایه خطی یکنواخت (\lr{ULA}) و آرایه مستطیلی یکنواخت (\lr{URA})، مورد بحث قرار گرفته و روابط مربوط به داده‌ها و ماتریس هدایت آرایه آن‌ها استخراج خواهد شد.



\subsection{آرایه خطی یکنواخت}
\label{subsec:3-3-1}

اکنون مورد آرایه خطی یکنواخت (\lr{ULA}) را در نظر بگیرید که از جمله رایج‌ترین آرایه‌های مورد استفاده در عمل است. یک نمونه پیکربندی در شکل \ref{fig:3-2} نشان داده شده است. همان‌طور که در بخش \ref{sec:3-2-1} توصیف شد، بردار هدایت آرایه متناظر با فرکانس مکانی $\mu_i = - \frac{2\pi}{\lambda} \Delta \sin \theta_i$ را می‌توان به‌صورت زیر نوشت:
\begin{equation}
	\bm{a}(\mu_i) = [1, e^{-j\mu_i}, e^{-j2\mu_i}, \dots, e^{-j(M-1)\mu_i}]^T, \quad 1 \leq i \leq d
	\label{eq:3-20}
\end{equation}
و ماتریس هدایت $\bm{A}$ دارای ساختار وندرموند\LTRfootnote{Vandermonde structure} زیر است:
\begin{equation}
	\bm{A} = [\bm{a}(\mu_1), \bm{a}(\mu_2), \dots, \bm{a}(\mu_d)] = 
	\begin{bmatrix} 
		1 & 1 & \dots & 1 \\
		e^{-j\mu_1} & e^{-j\mu_2} & \dots & e^{-j\mu_d} \\
		\vdots & \vdots & \ddots & \vdots \\
		e^{-j(M-1)\mu_1} & e^{-j(M-1)\mu_2} & \dots & e^{-j(M-1)\mu_d}
	\end{bmatrix}
	\label{eq:3-21}
\end{equation}
آرایه‌های \lr{ULA} مرکز-متقارن هستند، زیرا رابطه زیر به‌سادگی قابل اثبات است:
\begin{equation}
	\bm{\Pi}_M \bar{\bm{A}} = \bm{A} \bm{\Lambda}_\mu
	\label{eq:3-22}
\end{equation}
که در آن
\begin{equation}
	\bm{\Lambda}_\mu = 
	\begin{bmatrix} 
		e^{j(M-1)\mu_1} & 0 & \dots & 0 \\
		0 & e^{j(M-1)\mu_2} & \dots & 0 \\
		\vdots & \vdots & \ddots & \vdots \\
		0 & 0 & \dots & e^{j(M-1)\mu_d}
	\end{bmatrix}
	\nonumber
\end{equation}


بار دیگر، اگر مرکز آرایه به‌عنوان مرجع فاز\LTRfootnote{Phase reference} انتخاب شود، ماتریس هدایت آرایه حاصل، $\bm{A}_c$، را می‌توان با مقیاس‌دهی ستون‌های $\bm{A}$ به صورت زیر به‌دست آورد:
\begin{equation}
	\bm{A}_c = \bm{A} \cdot \bm{\Phi}^{-\frac{M-1}{2}}
	\label{eq:3-23}
\end{equation}
نشان دادن این موضوع دشوار نیست که $\bm{A}_c$ در رابطه \eqref{eq:3-13} صدق می‌کند؛ بنابراین $\bm{A}_c$ مرکز-متقارن است. این نتیجه در فصل \ref{chap:5} مورد استفاده قرار خواهد گرفت. بردارهای هدایت این آرایه، یعنی ستون‌های $\bm{A}_c$ عبارت‌اند از:
\begin{equation}
	\bm{a}_M(\mu_i) = e^{j\frac{M-1}{2}\mu_i} [1, e^{-j\mu_i}, e^{-j2\mu_i}, \dots, e^{-j(M-1)\mu_i}]^T, \quad 1 \leq i \leq d
	\label{eq:3-24}
\end{equation}
در اینجا، زیرنویس $M$ نشان‌دهنده بعد بردار هدایت حقیقی-چپ-$\bm{\Pi}$ یعنی $\bm{a}_M(\mu_i)$ است.

\subsection{آرایه مستطیلی یکنواخت (\lr{URA})}
\label{subsec:3-3-2}
زمانی که هدف، تخمین زوایای یک منبع در هر دو حالت سمت\LTRfootnote{Azimuth} و ارتفاع\LTRfootnote{Elevation} باشد، مطابق شکل \ref{fig:3-3} از آرایه‌های دوبعدی استفاده می‌شود. نمونه‌ای از پیکربندی چنین آرایه‌ای در شکل \ref{fig:3-4} نشان داده شده است. این شکل شامل سه فرم مختلف از جای‌گذاری عناصر آرایه است؛ برخی بدون عناصر مرکزی [شکل \ref{fig:3-4}(b)]، در حالی که برخی دیگر شامل خطوط متعامد از عناصر آرایه‌ای با مراکز فاز مشترک هستند [شکل \ref{fig:3-4}(a, c)] \cite{6}.

\begin{figure}[ht]
	\centering
	%\includegraphics[width=0.7\textwidth]{fig3-3}
	\caption{تعریف زوایای سمت و ارتفاع.}
	\label{fig:3-3}
\end{figure}

\begin{figure}[ht]
	\centering
	%\includegraphics[width=0.8\textwidth]{fig3-4}
	\caption{پیکربندی‌های آرایه دوبعدی مرکز-متقارن (الف-ج).}
	\label{fig:3-4}
\end{figure}


حال شکل \ref{fig:3-4}(الف) را در نظر می‌گیریم. ابعاد آرایه دوبعدی $M_x \times M_y$ است. عناصر آرایه با $(k_x, k_y)$ شماره‌گذاری می‌شوند که در آن $1 \leq k_x \leq M_x$ و $1 \leq k_y \leq M_y$ است. با پیروی از مدل داده‌ای که پیش‌تر توصیف شد، فرض می‌کنیم $d$ سیگنال باندباریک با طول موج $\lambda$ توسط $d$ منبع گسیل می‌شوند و جبهه موج تختِ سیگنال از منبع $i$-ام با زاویه سمت $\phi_i$ و زاویه ارتفاع $\theta_i$ بر آرایه می‌تابد ($1 \leq i \leq d$) (شکل \ref{fig:3-3} را ببینید). فرض کنید:
\begin{equation}
	u_i = \cos \phi_i \sin \theta_i, \quad v_i = \sin \phi_i \sin \theta_i, \quad 1 \leq i \leq d
	\label{eq:3-25}
\end{equation}
قرار می‌دهیم:
\begin{equation}
	\xi_i = \sin \theta_i e^{j \phi_i} = u_i + jv_i
	\label{eq:3-26}
\end{equation}
فرمول ساده‌ای برای تعیین زاویه سمت $\phi_i$ و ارتفاع $\theta_i$ به‌ترتیب از روی کسینوس‌های هادی\LTRfootnote{Direction cosines} متناظر $u_i$ و $v_i$ به‌صورت زیر ارائه می‌شود:
\begin{equation}
	\phi_i = \arg(\xi_i) \quad \text{و} \quad \theta_i = \arcsin(|\xi_i|)
	\label{eq:3-27}
\end{equation}
معادله \eqref{eq:3-27} نشان می‌دهد که تعیین \lr{DOA} یک منبع، معادل با تعیین کسینوس‌های هادی متناظر $u_i$ و $v_i$ است.

اگر داده‌های دریافت شده توسط عنصر آرایه $(k_x, k_y)$ در اسنپ‌شات\LTRfootnote{Snapshot} زمانی $t = t_n$ به‌صورت $x(t_n, k_x, k_y)$ داده شود، در هر اسنپ‌شات $M = M_x \cdot M_y$ نمونه وجود خواهد داشت. با تعمیم نتیجه یک‌بعدی \eqref{eq:3-9}، یافتن برهم‌نهی $d$ سیگنال و نویزهای دریافت شده توسط این عنصر دشوار نیست که به‌صورت زیر نوشته می‌شود:











